% Appendix A. The disclosed series in full, plus the exception comparison
% that §6 defers here. Tables generated by build/make_tables.py.

\section{The disclosed series in full}\label{sec:appendix_a}

Table~\ref{tab:manifest} identifies the ten filings every disclosed figure in
this paper is read from. Accession numbers are given so that any figure can be
traced to the document it came from and retrieved independently.

% Generated by build/make_tables.py from gs_filing_manifest.csv. Do not edit.

\begin{table}[tbp]
\centering
\small
\caption{The filings the disclosed series is read from}
\label{tab:manifest}
\begin{tabular}{@{}llll@{}}
\toprule
Fiscal year & Period end & Filed & Accession number \\
\midrule
FY2016 & 2016-12-31 & 2017-02-27 & \texttt{0001193125-17-056804} \\
FY2017 & 2017-12-31 & 2018-02-26 & \texttt{0001193125-18-056383} \\
FY2018 & 2018-12-31 & 2019-02-26 & \texttt{0001193125-19-050198} \\
FY2019 & 2019-12-31 & 2020-02-21 & \texttt{0001193125-20-043853} \\
FY2020 & 2020-12-31 & 2021-02-22 & \texttt{0001193125-21-049380} \\
FY2021 & 2021-12-31 & 2022-02-25 & \texttt{0001193125-22-052682} \\
FY2022 & 2022-12-31 & 2023-02-24 & \texttt{0000886982-23-000003} \\
FY2023 & 2023-12-31 & 2024-02-23 & \texttt{0000886982-24-000006} \\
FY2024 & 2024-12-31 & 2025-02-27 & \texttt{0000886982-25-000005} \\
FY2025 & 2025-12-31 & 2026-02-25 & \texttt{0000886982-26-000091} \\
\bottomrule
\end{tabular}
\notetext{The Goldman Sachs Group, Inc., Annual Report on Form 10-K, CIK 0000886982. Each is retrievable from EDGAR at \path{https://www.sec.gov/Archives/edgar/data/886982/<accession>/}. Every figure in this paper is taken from the Market Risk Management section of one of these ten documents.}
\end{table}


Table~\ref{tab:disclosed} in Section~\ref{sec:disclosed} gives the average
daily series, which is what the analysis uses. Table~\ref{tab:periodend} gives
the period-end series for completeness. The two measure different things: the
average daily figure summarises a year of risk-taking, while the period-end
figure is a single day's snapshot and is correspondingly noisier. The
diversification identity holds in both.

% Generated by build/make_tables.py from gs_var_period_end.csv. Do not edit.

\begin{table}[tbp]
\centering
\small
\caption{Period-end Value-at-Risk by risk category, FY2016--FY2025}
\label{tab:periodend}
\begin{tabular}{@{}lrrrrrrrr@{}}
\toprule
 & \multicolumn{4}{c}{Risk category} & & & & \\
\cmidrule(lr){2-5}
Year & Rates & Equities & FX & Commod. & Summed & Diversif. & Total & \% of sum \\
\midrule
2016 & 45 & 34 & 23 & 19 & 121 & -39 & 82 & 32.2 \\
2017 & 48 & 31 & 7 & 9 & 95 & -30 & 65 & 31.6 \\
2018 & 46 & 32 & 12 & 11 & 101 & -44 & 57 & 43.6 \\
2019 & 54 & 24 & 10 & 10 & 98 & -28 & 70 & 28.6 \\
2020 & 60 & 50 & 11 & 16 & 137 & -46 & 91 & 33.6 \\
2021 & 69 & 31 & 19 & 30 & 149 & -58 & 91 & 38.9 \\
2022 & 108 & 27 & 35 & 18 & 188 & -85 & 103 & 45.2 \\
2023 & 93 & 25 & 15 & 14 & 147 & -54 & 93 & 36.7 \\
2024 & 84 & 41 & 38 & 14 & 177 & -86 & 91 & 48.6 \\
2025 & 59 & 45 & 17 & 18 & 139 & -60 & 79 & 43.2 \\
\bottomrule
\end{tabular}
\notetext{In millions of US dollars, except the final column. The analysis in this paper uses average daily VaR (Table~\ref{tab:disclosed}); the period-end series is a single-day snapshot and is reported here for completeness. The identity holds in this series too.}
\end{table}


Section~\ref{sec:results} uses the period-end series to separate the amplitude
of the disclosed variation from its persistence, and Table~\ref{tab:amplitude}
carries that comparison. The point in brief: measured period-end Goldman's own
ratio has a year-to-year standard deviation of 6.72 percentage points against
2.08 for the average daily series, so the disclosed steadiness is a property of
the annual mean rather than of the ratio itself.

The disclosed high and low VaR by category point the same way. Taken as first
reported, the median within-year ratio of high to low runs from 1.9 for
interest rates to 4.3 for currency rates. A book whose category VaRs move
through ranges that wide over a year is not a book whose diversification effect
sits still from day to day, and the filings are explicit enough about this that
no inference is needed.

\subsection*{Exception counts}

Section~\ref{sec:results} defers this comparison here because it cannot carry
weight. Table~\ref{tab:exceptions} gives it.

% Generated by build/make_tables.py from exceptions_vs_disclosed.csv. Do not edit.

\begin{table}[tbp]
\centering
\small
\caption{Exception counts at 95\%: the proxy book against Goldman's disclosed counts}
\label{tab:exceptions}
\begin{tabular}{@{}lrrrr@{}}
\toprule
 & \multicolumn{3}{c}{Proxy book, plain historical simulation} & Goldman \\
\cmidrule(lr){2-4}
Year & Days & Exceptions & Rate (\%) & Exceptions \\
\midrule
2016 & 233 & 9 & 3.9 & 0 \\
2017 & 232 & 5 & 2.2 & 0 \\
2018 & 230 & 11 & 4.8 & 2 \\
2019 & 225 & 8 & 3.6 & 0 \\
2020 & 227 & 12 & 5.3 & 2 \\
2021 & 231 & 11 & 4.8 & 1 \\
2022 & 230 & 70 & 30.4 & 2 \\
2023 & 233 & 19 & 8.2 & 1 \\
2024 & 228 & 11 & 4.8 & 2 \\
2025 & 224 & 4 & 1.8 & 0 \\
\midrule
Total & 2,293 & 160 & 6.98 & 10 \\
\bottomrule
\end{tabular}
\notetext{\textbf{These two columns do not measure the same thing and should not be compared as though they did.} Goldman scores VaR against daily net revenues, which include intraday market-making income that is positive far more often than not; the proxy book has no client flow and no bid--offer spread, so its profit and loss is a clean mark on a static position. The table is reported for completeness and the gap is an indication of what the revenue cushion is worth, not evidence that either model is miscalibrated. No conclusion in this paper rests on it.}
\end{table}


Goldman recorded 10 exceptions at the 95\% risk-management measure across the
ten years, against 160 for the proxy book. The gap is not evidence about either
model's calibration, because the two are scored against different quantities:
Goldman compares VaR against daily net revenues, which include intraday
market-making income that is positive on the large majority of days, while the
proxy book is a static position marked cleanly with no client flow and no
bid--offer spread. The filings make this point themselves, observing that in
periods with far more positive than negative net revenue days, fewer exceptions
are expected than the confidence level implies.

The distance between 0.40\% and 6.98\% of days is therefore best read as an
indication of how much the revenue cushion is worth, and it is large. It is not
used anywhere in this paper's conclusions.
